\chapter{Conclusão}

\indent Neste trabalho foi realizada a investigação do método dos elementos virtuais no estado plano de tensões e deformações em problemas de vibração livre. Para tal, foram calculados três operadores de projeção com formulação genérica possibilitando a utilização de aproximações lineares e de alta ordem. A formulação foi aplicada a três problemas de vibração livre: uma placa engastada com um furo no centro, uma barragem de terra e uma parede estrutural. Esses problemas foram comparados com os resultados obtidos pelo Método dos Elementos Finitos e o Método dos Elementos Finitos Generalizados.\\
\indent Avaliando o primeiro problema foram sugeridas quatro malhas iniciais com aumento progressivo da quantidade de nós na fronteira com o furo onde foi possível observar a aproximação da solução numérica à de referência com o aumento dos graus de liberdade. Em uma tentativa de reduzir o erro induzido pelas funções não polinomiais na região da fronteira com o furo estabeleceu-se um novo conjunto de malhas aumentando a quantidade de elementos no domínio.\\
\indent O aumento de elementos resultou em erros maiores a partir da quarta frequência. Investigando-se o condicionamento, verificou-se um aumento do número de condicionamento relacionado à quantidade de elementos. Conclui-se que as funções ao redor do furo estavam sendo completamente capturadas pelo termo de estabilização. A estabilização é escalada por um parâmetro que deve ser calibrado para o problema. O aumento da quantidade de elementos na região da fronteira com o furo, aliado à diminuição dos graus de liberdade que representavam a geometria circular, acarretou o aumento dos erros da parcela de estabilização que por consequência afetam os resultados em malhas com maior quantidade elementos. \\
\indent No estudo do ganho de precisão em comparação ao MEF, para a malha com apenas 34 graus de liberdade houve uma perda de precisão que logo foi superada com a adição de novos graus de liberdade. Esse aumento só é válido para a malha com 8 elementos, o aumento do erro provindo da estabilização torna as malhas de 34 e 64 elementos piores que as do MEF. Comparando os resultados obtidos pelo MEV e o MEFG nota-se que em todos os casos o MEV atingiu melhores resultados que o MEF e o MEFG para a malha de 8 elementos, com o menor erro sendo equivalente ao do MEFG para a malha analisada.\\
\indent Mesmo com resultados melhores que os dois métodos, o erro do MEV equivale a $\approx 10\%$ o que torna necessário futuras pesquisas para melhorar à resposta. Aqui se sugerem estudos de envolvendo elementos com formulação para arestas curvas, a aplicação da formulação enriquecida de elementos virtuais e uma comparação de resultados com a formulação sem estabilização do MEV.\\
\indent Para o problema da barragem de terra tentou-se aproximar a solução utilizando apenas 3 elementos com aumento da quantidade de nós entre uma malha e outra. Das 3 malhas utilizadas somente a primeira frequência natural não apresenta modos de vibração artificias com destaque para as malhas com adições que graus de liberdade que obtiveram erros muito maior que a malha mais grosseira. Após revisão da literatura concluiu-se que o erro acontece porque com a adição de nós a distância entre eles diminuia tornando a matriz $\mathcal{D}$ linearmente dependente afetando a estabilização do método. Ao utilizar uma malha com mais elementos e uma melhor distribuição dos graus de liberdade se obteve resultados mais próximos da solução de referência. \\
\indent Para melhor aproximar a solução foi testada a aproximação quadrática. Para as malhas com adições de graus de liberdade nas arestas os resultados obtidos foram os mesmos para aproximação $k=1$. Todavia, para a malha mais grosseira e a malha equilibrada obtiveram resultados muito melhores que sua contraparte de baixa ordem, com destaque para a malha equilibrada que obteve resultado quase idêntico à solução de referência. Comparando os resultados entre a malha mais grosseira e a mais refinada, percebe-se que apenas as malhas equilibradas possuem melhora de discretização, demonstrando necessidade de equilíbrio entre a quantidade de elementos e de graus de liberdade para obtenção de uma solução precisa. Comparando ao resultado obtido pelo MEF percebe-se que apenas as malhas equilibradas possuem resultados melhores que o MEF, e entre as duas malhas, apenas a malha quadrática obteve erros inferiores ao MEFG.\\
\indent No problema da parede estrutural foram testadas 4 malhas de primeira ordem com aumento da quantidade de elementos. Todas as malhas obtiveram altos erros. A análise indicou fenômeno de travamento por cisalhamento, levando a adoção de malhas de segunda ordem para a continuação da investigação. Mesmo tendo sido um sucesso, para a diminuição do erro substitui-se o parâmetro de estabilização para um com a parte elástica baseado unicamente nas propriedades do material e a parte de inércia variando com o diâmetro do elemento. Notou-se que essa mudança diminuia o erro da solução.\\
\indent Continuando as investigações, tomou-se o parâmetro de estabilização de inércia como nulo. Para malha mais refinada não havia mudanças, porém para a malha mais grosseira seu resultado foi quase o mesmo da malha mais refinada, indicando que os erros da estabilização da massa afetavam negativamente e para os problemas de vibração livre não amortecida a parcela de inércia não precisa ser estabilizada, apenas a parte elástica. Mesmo com essas conclusões, determinou-se que grande parte do erro originava-se na escolha da base polinomial, de maneira que o aumento do grau polinomial faz com que a base perca sua linearidade, afetando ambas as parcelas de consistência e estabilização. \\
\indent Considerando os resultados obtidos, são sugeridos estudos do efeito de diferentes bases polinomiais, assim como o efeito das diferentes formas de estabilização para a rigidez de modo a evitar o enrijecimento artificial causado pelos erros de estabilização.